\section{Contexto e relação com trabalho anterior}
Sparse MoE e conditional computation antecedem o Tesy \cite{shazeer2017moe,fedus2022switch}; o próprio gpt-oss é um Transformer MoE \cite{openai2025gptoss}. Offload multi-tier para LLMs já foi estudado em FlexGen e \emph{LLM in a flash} \cite{sheng2023flexgen,alizadeh2024llmflash}. Para MoE, MoE-Infinity explora tracing, cache e prefetch de experts \cite{xue2024moeinfinity}; Fiddler e KTransformers cobrem orchestration CPU--GPU \cite{kamahori2025fiddler,chen2025ktransformers}; FineMoE estuda offload fine-grained \cite{yu2026finemoe}; e ExpertFlow combina previsão, caching e token scheduling \cite{he2026expertflow}.

O espaço contemporâneo aproxima-se ainda mais do problema Tesy. Wang et al. apresentam inferência MoE local CPU--GPU com stream-loading e overlap \cite{wang2026localslo}. Kairox usa balancing online, predictive prefetch e uma cache orientada à localidade de ativações em PCs de consumo \cite{jiang2026kairox}. CoPilotIO mostra que o caminho storage--GPU e os stalls de I/O são, eles próprios, objeto de co-design CPU/GPU \cite{chen2026copilotio}. \emph{Patterns Behind Chaos}, publicado em ISCA'26, estuda traces de seleção de experts e usa-os para orientar placement \cite{yu2026patterns}.

Há também prior art particularmente próximo do contrato de exactness: SpecPrefetch usa um predictor apenas para transfer antecipado, mantendo o router nativo congelado como decisor dos experts efetivamente executados \cite{kong2026specprefetch}. Logo, “prediction sem alterar routing” é uma restrição correta do Tesy, mas não uma claim de novidade.

Os resultados externos não são projetados para o Ryzen/RTX do Tesy. Diferenças de modelo, precisão, hardware, contexto, batch, storage e definição de throughput impedem transferir speedups da literatura para este target.
